% page 295 du fac-simile (image p-294 du scan)
% bloc 160.0 x 237.7 mm ; marge gauche 42.0 mm
\pgc{20.320mm}{0.000mm}{
\l{\cel{54}287}
\l{}
\l{}
\l{kolportar. (\sou{trans.}) Portar (vari) aden diversa loki, regioni,}
\l{\cel{3}e c., (por disvendar li en ca loki, regioni). - DEF.}
\l{}
\l{kolubro. (\sou{zool.}) Reptero ne-venenoza, ek la familio \textquotedbl{}serpenti\textquotedbl{}.}
\l{\cel{3}- L.\cel{1}\sou{coluber}. - EFSL.}
\l{}
\l{\sou{kolumbario}. (en\cel{1}\sou{Roma, antique}). Konstrukturo funerala qua kon-}
\l{\cel{3}tenis nichi aden quionu pozis la urni mortintala. - DEFIS.}
\l{}
\l{kolumno. Dividuro vertikala (kolonatra), inter-paralela, di}
\l{\cel{3}la paginedi de komposturo, de imprimuro (texto). - DE.}
\l{}
\l{koluro. (\sou{astron.})\cel{2}Singla de la du granda cirkli di la Tero-}
\l{\cel{3}sfero, qui intersekas ortangule ye la poli, e pasas, ica ye}
\l{\cel{3}la punti equinoxala, ita ye la punti solsticala. - DEFIS.}
\l{}
\l{koluteo. (\sou{bot.)}\cel{2}Arboreto, ek la familio \textquotedbl{}papilionacei\textquotedbl{}, di qua}
\l{\cel{3}la folii olim uzesis kom purgivo, e di qua la sheli, pro}
\l{\cel{3}explozar bruisoze, esas amuzivo por la infanti. - L.\cel{1}\sou{colutea}}
\l{\cel{3}\sou{arborescens}.- IL.}
\l{}
\l{koluzionar. (\sou{netrans}.). (\sou{aludante du personi qui su prizentas}}
\l{\cel{3}a la judiciisto). Interkonvencionar sekrete, detrimente di}
\l{\cel{3}triesma koncernato di la proceso. - DEFIS.}
\l{}
\l{kolzo. (\sou{bot.}) Planto krucifera, di qua la grano donas oleo}
\l{\cel{3}quan onu uzas precipue por la lumizado. - L.\cel{1}\sou{brassica napus}}
\l{\cel{3}\sou{oleifera}.- EFIRS.}
\l{}
\l{kom.\cel{3}\textquotedbl{}Havante la qualeso, la titulo, la ofico, la rolo, di\textquotedbl{}.}
\l{\cel{3}- F.}
\l{}
\l{komo.\cel{2}+Puntuo-signo (,) qua indikas ke oportas (dicionale)}
\l{\cel{3}separar ula membri di la frazo, sen haltar longe. - DEF.}
\l{}
\l{komao. (\sou{muziko}). Intervalo muzikala tre mikra, quan onu}
\l{\cel{3}desegardas, ordinare - pro ke ol esas apene orelo-percepte-}
\l{\cel{3}bla. - DEF.}
\l{}
\l{komandar.\cel{2}(\sou{trans}.)\cel{2}(\sou{aludante armeo-chefo)}. Signalar ke movo,}
\l{\cel{3}manovro, esas agenda da la subordinito - impero extreme}
\l{\cel{3}severa, rigoroza ed absoluta, qua supresas omna responso che}
\l{\cel{3}la imperario. - II. (tekn.) (\sou{aludante organo di mekanismo)}.}
\l{\cel{3}Komunikar sua movo ad altra organo per ingranajo, pulio,}
\l{\cel{3}rimeno, e c.) DEFIRS.}
\l{}
\l{komanditar. (\sou{trans.)}\cel{3}Subtenar (entraprezajo o la persono qua}
\l{\cel{3}entraprezas) kom nura pekunio-prestero, sen intervenar pri}
\l{\cel{3}quale uzesas ica pekunio. - DEFIS.}
\l{}
\l{\cel{1}\sou{komando}ro. La persono qua esas super la oficiro, en la ordeni}
\l{\cel{3}honoriziva di la +chevalieraro. - DEFIRS.}
}
